\documentclass[11pt]{article}

    \usepackage[breakable]{tcolorbox}
    \usepackage{parskip} % Stop auto-indenting (to mimic markdown behaviour)
    

    % Basic figure setup, for now with no caption control since it's done
    % automatically by Pandoc (which extracts ![](path) syntax from Markdown).
    \usepackage{graphicx}
    % Keep aspect ratio if custom image width or height is specified
    \setkeys{Gin}{keepaspectratio}
    % Maintain compatibility with old templates. Remove in nbconvert 6.0
    \let\Oldincludegraphics\includegraphics
    % Ensure that by default, figures have no caption (until we provide a
    % proper Figure object with a Caption API and a way to capture that
    % in the conversion process - todo).
    \usepackage{caption}
    \DeclareCaptionFormat{nocaption}{}
    \captionsetup{format=nocaption,aboveskip=0pt,belowskip=0pt}

    \usepackage{float}
    \floatplacement{figure}{H} % forces figures to be placed at the correct location
    \usepackage{xcolor} % Allow colors to be defined
    \usepackage{enumerate} % Needed for markdown enumerations to work
    \usepackage{geometry} % Used to adjust the document margins
    \usepackage{amsmath} % Equations
    \usepackage{amssymb} % Equations
    \usepackage{textcomp} % defines textquotesingle
    % Hack from http://tex.stackexchange.com/a/47451/13684:
    \AtBeginDocument{%
        \def\PYZsq{\textquotesingle}% Upright quotes in Pygmentized code
    }
    \usepackage{upquote} % Upright quotes for verbatim code
    \usepackage{eurosym} % defines \euro

    \usepackage{iftex}
    \ifPDFTeX
        \usepackage[T1]{fontenc}
        \IfFileExists{alphabeta.sty}{
              \usepackage{alphabeta}
          }{
              \usepackage[mathletters]{ucs}
              \usepackage[utf8x]{inputenc}
          }
    \else
        \usepackage{fontspec}
        \usepackage{unicode-math}
    \fi

    \usepackage{fancyvrb} % verbatim replacement that allows latex
    \usepackage{grffile} % extends the file name processing of package graphics
                         % to support a larger range
    \makeatletter % fix for old versions of grffile with XeLaTeX
    \@ifpackagelater{grffile}{2019/11/01}
    {
      % Do nothing on new versions
    }
    {
      \def\Gread@@xetex#1{%
        \IfFileExists{"\Gin@base".bb}%
        {\Gread@eps{\Gin@base.bb}}%
        {\Gread@@xetex@aux#1}%
      }
    }
    \makeatother
    \usepackage[Export]{adjustbox} % Used to constrain images to a maximum size
    \adjustboxset{max size={0.9\linewidth}{0.9\paperheight}}

    % The hyperref package gives us a pdf with properly built
    % internal navigation ('pdf bookmarks' for the table of contents,
    % internal cross-reference links, web links for URLs, etc.)
    \usepackage{hyperref}
    % The default LaTeX title has an obnoxious amount of whitespace. By default,
    % titling removes some of it. It also provides customization options.
    \usepackage{titling}
    \usepackage{longtable} % longtable support required by pandoc >1.10
    \usepackage{booktabs}  % table support for pandoc > 1.12.2
    \usepackage{array}     % table support for pandoc >= 2.11.3
    \usepackage{calc}      % table minipage width calculation for pandoc >= 2.11.1
    \usepackage[inline]{enumitem} % IRkernel/repr support (it uses the enumerate* environment)
    \usepackage[normalem]{ulem} % ulem is needed to support strikethroughs (\sout)
                                % normalem makes italics be italics, not underlines
    \usepackage{soul}      % strikethrough (\st) support for pandoc >= 3.0.0
    \usepackage{mathrsfs}
    

    
    % Colors for the hyperref package
    \definecolor{urlcolor}{rgb}{0,.145,.698}
    \definecolor{linkcolor}{rgb}{.71,0.21,0.01}
    \definecolor{citecolor}{rgb}{.12,.54,.11}

    % ANSI colors
    \definecolor{ansi-black}{HTML}{3E424D}
    \definecolor{ansi-black-intense}{HTML}{282C36}
    \definecolor{ansi-red}{HTML}{E75C58}
    \definecolor{ansi-red-intense}{HTML}{B22B31}
    \definecolor{ansi-green}{HTML}{00A250}
    \definecolor{ansi-green-intense}{HTML}{007427}
    \definecolor{ansi-yellow}{HTML}{DDB62B}
    \definecolor{ansi-yellow-intense}{HTML}{B27D12}
    \definecolor{ansi-blue}{HTML}{208FFB}
    \definecolor{ansi-blue-intense}{HTML}{0065CA}
    \definecolor{ansi-magenta}{HTML}{D160C4}
    \definecolor{ansi-magenta-intense}{HTML}{A03196}
    \definecolor{ansi-cyan}{HTML}{60C6C8}
    \definecolor{ansi-cyan-intense}{HTML}{258F8F}
    \definecolor{ansi-white}{HTML}{C5C1B4}
    \definecolor{ansi-white-intense}{HTML}{A1A6B2}
    \definecolor{ansi-default-inverse-fg}{HTML}{FFFFFF}
    \definecolor{ansi-default-inverse-bg}{HTML}{000000}

    % common color for the border for error outputs.
    \definecolor{outerrorbackground}{HTML}{FFDFDF}

    % commands and environments needed by pandoc snippets
    % extracted from the output of `pandoc -s`
    \providecommand{\tightlist}{%
      \setlength{\itemsep}{0pt}\setlength{\parskip}{0pt}}
    \DefineVerbatimEnvironment{Highlighting}{Verbatim}{commandchars=\\\{\}}
    % Add ',fontsize=\small' for more characters per line
    \newenvironment{Shaded}{}{}
    \newcommand{\KeywordTok}[1]{\textcolor[rgb]{0.00,0.44,0.13}{\textbf{{#1}}}}
    \newcommand{\DataTypeTok}[1]{\textcolor[rgb]{0.56,0.13,0.00}{{#1}}}
    \newcommand{\DecValTok}[1]{\textcolor[rgb]{0.25,0.63,0.44}{{#1}}}
    \newcommand{\BaseNTok}[1]{\textcolor[rgb]{0.25,0.63,0.44}{{#1}}}
    \newcommand{\FloatTok}[1]{\textcolor[rgb]{0.25,0.63,0.44}{{#1}}}
    \newcommand{\CharTok}[1]{\textcolor[rgb]{0.25,0.44,0.63}{{#1}}}
    \newcommand{\StringTok}[1]{\textcolor[rgb]{0.25,0.44,0.63}{{#1}}}
    \newcommand{\CommentTok}[1]{\textcolor[rgb]{0.38,0.63,0.69}{\textit{{#1}}}}
    \newcommand{\OtherTok}[1]{\textcolor[rgb]{0.00,0.44,0.13}{{#1}}}
    \newcommand{\AlertTok}[1]{\textcolor[rgb]{1.00,0.00,0.00}{\textbf{{#1}}}}
    \newcommand{\FunctionTok}[1]{\textcolor[rgb]{0.02,0.16,0.49}{{#1}}}
    \newcommand{\RegionMarkerTok}[1]{{#1}}
    \newcommand{\ErrorTok}[1]{\textcolor[rgb]{1.00,0.00,0.00}{\textbf{{#1}}}}
    \newcommand{\NormalTok}[1]{{#1}}

    % Additional commands for more recent versions of Pandoc
    \newcommand{\ConstantTok}[1]{\textcolor[rgb]{0.53,0.00,0.00}{{#1}}}
    \newcommand{\SpecialCharTok}[1]{\textcolor[rgb]{0.25,0.44,0.63}{{#1}}}
    \newcommand{\VerbatimStringTok}[1]{\textcolor[rgb]{0.25,0.44,0.63}{{#1}}}
    \newcommand{\SpecialStringTok}[1]{\textcolor[rgb]{0.73,0.40,0.53}{{#1}}}
    \newcommand{\ImportTok}[1]{{#1}}
    \newcommand{\DocumentationTok}[1]{\textcolor[rgb]{0.73,0.13,0.13}{\textit{{#1}}}}
    \newcommand{\AnnotationTok}[1]{\textcolor[rgb]{0.38,0.63,0.69}{\textbf{\textit{{#1}}}}}
    \newcommand{\CommentVarTok}[1]{\textcolor[rgb]{0.38,0.63,0.69}{\textbf{\textit{{#1}}}}}
    \newcommand{\VariableTok}[1]{\textcolor[rgb]{0.10,0.09,0.49}{{#1}}}
    \newcommand{\ControlFlowTok}[1]{\textcolor[rgb]{0.00,0.44,0.13}{\textbf{{#1}}}}
    \newcommand{\OperatorTok}[1]{\textcolor[rgb]{0.40,0.40,0.40}{{#1}}}
    \newcommand{\BuiltInTok}[1]{{#1}}
    \newcommand{\ExtensionTok}[1]{{#1}}
    \newcommand{\PreprocessorTok}[1]{\textcolor[rgb]{0.74,0.48,0.00}{{#1}}}
    \newcommand{\AttributeTok}[1]{\textcolor[rgb]{0.49,0.56,0.16}{{#1}}}
    \newcommand{\InformationTok}[1]{\textcolor[rgb]{0.38,0.63,0.69}{\textbf{\textit{{#1}}}}}
    \newcommand{\WarningTok}[1]{\textcolor[rgb]{0.38,0.63,0.69}{\textbf{\textit{{#1}}}}}
    \makeatletter
    \newsavebox\pandoc@box
    \newcommand*\pandocbounded[1]{%
      \sbox\pandoc@box{#1}%
      % scaling factors for width and height
      \Gscale@div\@tempa\textheight{\dimexpr\ht\pandoc@box+\dp\pandoc@box\relax}%
      \Gscale@div\@tempb\linewidth{\wd\pandoc@box}%
      % select the smaller of both
      \ifdim\@tempb\p@<\@tempa\p@
        \let\@tempa\@tempb
      \fi
      % scaling accordingly (\@tempa < 1)
      \ifdim\@tempa\p@<\p@
        \scalebox{\@tempa}{\usebox\pandoc@box}%
      % scaling not needed, use as it is
      \else
        \usebox{\pandoc@box}%
      \fi
    }
    \makeatother

    % Define a nice break command that doesn't care if a line doesn't already
    % exist.
    \def\br{\hspace*{\fill} \\* }
    % Math Jax compatibility definitions
    \def\gt{>}
    \def\lt{<}
    \let\Oldtex\TeX
    \let\Oldlatex\LaTeX
    \renewcommand{\TeX}{\textrm{\Oldtex}}
    \renewcommand{\LaTeX}{\textrm{\Oldlatex}}
    % Document parameters
    % Document title
    \title{Modelos y Sistemas - Trabajo Final}
    
    
    
    
    
    
    
% Pygments definitions
\makeatletter
\def\PY@reset{\let\PY@it=\relax \let\PY@bf=\relax%
    \let\PY@ul=\relax \let\PY@tc=\relax%
    \let\PY@bc=\relax \let\PY@ff=\relax}
\def\PY@tok#1{\csname PY@tok@#1\endcsname}
\def\PY@toks#1+{\ifx\relax#1\empty\else%
    \PY@tok{#1}\expandafter\PY@toks\fi}
\def\PY@do#1{\PY@bc{\PY@tc{\PY@ul{%
    \PY@it{\PY@bf{\PY@ff{#1}}}}}}}
\def\PY#1#2{\PY@reset\PY@toks#1+\relax+\PY@do{#2}}

\@namedef{PY@tok@w}{\def\PY@tc##1{\textcolor[rgb]{0.73,0.73,0.73}{##1}}}
\@namedef{PY@tok@c}{\let\PY@it=\textit\def\PY@tc##1{\textcolor[rgb]{0.24,0.48,0.48}{##1}}}
\@namedef{PY@tok@cp}{\def\PY@tc##1{\textcolor[rgb]{0.61,0.40,0.00}{##1}}}
\@namedef{PY@tok@k}{\let\PY@bf=\textbf\def\PY@tc##1{\textcolor[rgb]{0.00,0.50,0.00}{##1}}}
\@namedef{PY@tok@kp}{\def\PY@tc##1{\textcolor[rgb]{0.00,0.50,0.00}{##1}}}
\@namedef{PY@tok@kt}{\def\PY@tc##1{\textcolor[rgb]{0.69,0.00,0.25}{##1}}}
\@namedef{PY@tok@o}{\def\PY@tc##1{\textcolor[rgb]{0.40,0.40,0.40}{##1}}}
\@namedef{PY@tok@ow}{\let\PY@bf=\textbf\def\PY@tc##1{\textcolor[rgb]{0.67,0.13,1.00}{##1}}}
\@namedef{PY@tok@nb}{\def\PY@tc##1{\textcolor[rgb]{0.00,0.50,0.00}{##1}}}
\@namedef{PY@tok@nf}{\def\PY@tc##1{\textcolor[rgb]{0.00,0.00,1.00}{##1}}}
\@namedef{PY@tok@nc}{\let\PY@bf=\textbf\def\PY@tc##1{\textcolor[rgb]{0.00,0.00,1.00}{##1}}}
\@namedef{PY@tok@nn}{\let\PY@bf=\textbf\def\PY@tc##1{\textcolor[rgb]{0.00,0.00,1.00}{##1}}}
\@namedef{PY@tok@ne}{\let\PY@bf=\textbf\def\PY@tc##1{\textcolor[rgb]{0.80,0.25,0.22}{##1}}}
\@namedef{PY@tok@nv}{\def\PY@tc##1{\textcolor[rgb]{0.10,0.09,0.49}{##1}}}
\@namedef{PY@tok@no}{\def\PY@tc##1{\textcolor[rgb]{0.53,0.00,0.00}{##1}}}
\@namedef{PY@tok@nl}{\def\PY@tc##1{\textcolor[rgb]{0.46,0.46,0.00}{##1}}}
\@namedef{PY@tok@ni}{\let\PY@bf=\textbf\def\PY@tc##1{\textcolor[rgb]{0.44,0.44,0.44}{##1}}}
\@namedef{PY@tok@na}{\def\PY@tc##1{\textcolor[rgb]{0.41,0.47,0.13}{##1}}}
\@namedef{PY@tok@nt}{\let\PY@bf=\textbf\def\PY@tc##1{\textcolor[rgb]{0.00,0.50,0.00}{##1}}}
\@namedef{PY@tok@nd}{\def\PY@tc##1{\textcolor[rgb]{0.67,0.13,1.00}{##1}}}
\@namedef{PY@tok@s}{\def\PY@tc##1{\textcolor[rgb]{0.73,0.13,0.13}{##1}}}
\@namedef{PY@tok@sd}{\let\PY@it=\textit\def\PY@tc##1{\textcolor[rgb]{0.73,0.13,0.13}{##1}}}
\@namedef{PY@tok@si}{\let\PY@bf=\textbf\def\PY@tc##1{\textcolor[rgb]{0.64,0.35,0.47}{##1}}}
\@namedef{PY@tok@se}{\let\PY@bf=\textbf\def\PY@tc##1{\textcolor[rgb]{0.67,0.36,0.12}{##1}}}
\@namedef{PY@tok@sr}{\def\PY@tc##1{\textcolor[rgb]{0.64,0.35,0.47}{##1}}}
\@namedef{PY@tok@ss}{\def\PY@tc##1{\textcolor[rgb]{0.10,0.09,0.49}{##1}}}
\@namedef{PY@tok@sx}{\def\PY@tc##1{\textcolor[rgb]{0.00,0.50,0.00}{##1}}}
\@namedef{PY@tok@m}{\def\PY@tc##1{\textcolor[rgb]{0.40,0.40,0.40}{##1}}}
\@namedef{PY@tok@gh}{\let\PY@bf=\textbf\def\PY@tc##1{\textcolor[rgb]{0.00,0.00,0.50}{##1}}}
\@namedef{PY@tok@gu}{\let\PY@bf=\textbf\def\PY@tc##1{\textcolor[rgb]{0.50,0.00,0.50}{##1}}}
\@namedef{PY@tok@gd}{\def\PY@tc##1{\textcolor[rgb]{0.63,0.00,0.00}{##1}}}
\@namedef{PY@tok@gi}{\def\PY@tc##1{\textcolor[rgb]{0.00,0.52,0.00}{##1}}}
\@namedef{PY@tok@gr}{\def\PY@tc##1{\textcolor[rgb]{0.89,0.00,0.00}{##1}}}
\@namedef{PY@tok@ge}{\let\PY@it=\textit}
\@namedef{PY@tok@gs}{\let\PY@bf=\textbf}
\@namedef{PY@tok@ges}{\let\PY@bf=\textbf\let\PY@it=\textit}
\@namedef{PY@tok@gp}{\let\PY@bf=\textbf\def\PY@tc##1{\textcolor[rgb]{0.00,0.00,0.50}{##1}}}
\@namedef{PY@tok@go}{\def\PY@tc##1{\textcolor[rgb]{0.44,0.44,0.44}{##1}}}
\@namedef{PY@tok@gt}{\def\PY@tc##1{\textcolor[rgb]{0.00,0.27,0.87}{##1}}}
\@namedef{PY@tok@err}{\def\PY@bc##1{{\setlength{\fboxsep}{\string -\fboxrule}\fcolorbox[rgb]{1.00,0.00,0.00}{1,1,1}{\strut ##1}}}}
\@namedef{PY@tok@kc}{\let\PY@bf=\textbf\def\PY@tc##1{\textcolor[rgb]{0.00,0.50,0.00}{##1}}}
\@namedef{PY@tok@kd}{\let\PY@bf=\textbf\def\PY@tc##1{\textcolor[rgb]{0.00,0.50,0.00}{##1}}}
\@namedef{PY@tok@kn}{\let\PY@bf=\textbf\def\PY@tc##1{\textcolor[rgb]{0.00,0.50,0.00}{##1}}}
\@namedef{PY@tok@kr}{\let\PY@bf=\textbf\def\PY@tc##1{\textcolor[rgb]{0.00,0.50,0.00}{##1}}}
\@namedef{PY@tok@bp}{\def\PY@tc##1{\textcolor[rgb]{0.00,0.50,0.00}{##1}}}
\@namedef{PY@tok@fm}{\def\PY@tc##1{\textcolor[rgb]{0.00,0.00,1.00}{##1}}}
\@namedef{PY@tok@vc}{\def\PY@tc##1{\textcolor[rgb]{0.10,0.09,0.49}{##1}}}
\@namedef{PY@tok@vg}{\def\PY@tc##1{\textcolor[rgb]{0.10,0.09,0.49}{##1}}}
\@namedef{PY@tok@vi}{\def\PY@tc##1{\textcolor[rgb]{0.10,0.09,0.49}{##1}}}
\@namedef{PY@tok@vm}{\def\PY@tc##1{\textcolor[rgb]{0.10,0.09,0.49}{##1}}}
\@namedef{PY@tok@sa}{\def\PY@tc##1{\textcolor[rgb]{0.73,0.13,0.13}{##1}}}
\@namedef{PY@tok@sb}{\def\PY@tc##1{\textcolor[rgb]{0.73,0.13,0.13}{##1}}}
\@namedef{PY@tok@sc}{\def\PY@tc##1{\textcolor[rgb]{0.73,0.13,0.13}{##1}}}
\@namedef{PY@tok@dl}{\def\PY@tc##1{\textcolor[rgb]{0.73,0.13,0.13}{##1}}}
\@namedef{PY@tok@s2}{\def\PY@tc##1{\textcolor[rgb]{0.73,0.13,0.13}{##1}}}
\@namedef{PY@tok@sh}{\def\PY@tc##1{\textcolor[rgb]{0.73,0.13,0.13}{##1}}}
\@namedef{PY@tok@s1}{\def\PY@tc##1{\textcolor[rgb]{0.73,0.13,0.13}{##1}}}
\@namedef{PY@tok@mb}{\def\PY@tc##1{\textcolor[rgb]{0.40,0.40,0.40}{##1}}}
\@namedef{PY@tok@mf}{\def\PY@tc##1{\textcolor[rgb]{0.40,0.40,0.40}{##1}}}
\@namedef{PY@tok@mh}{\def\PY@tc##1{\textcolor[rgb]{0.40,0.40,0.40}{##1}}}
\@namedef{PY@tok@mi}{\def\PY@tc##1{\textcolor[rgb]{0.40,0.40,0.40}{##1}}}
\@namedef{PY@tok@il}{\def\PY@tc##1{\textcolor[rgb]{0.40,0.40,0.40}{##1}}}
\@namedef{PY@tok@mo}{\def\PY@tc##1{\textcolor[rgb]{0.40,0.40,0.40}{##1}}}
\@namedef{PY@tok@ch}{\let\PY@it=\textit\def\PY@tc##1{\textcolor[rgb]{0.24,0.48,0.48}{##1}}}
\@namedef{PY@tok@cm}{\let\PY@it=\textit\def\PY@tc##1{\textcolor[rgb]{0.24,0.48,0.48}{##1}}}
\@namedef{PY@tok@cpf}{\let\PY@it=\textit\def\PY@tc##1{\textcolor[rgb]{0.24,0.48,0.48}{##1}}}
\@namedef{PY@tok@c1}{\let\PY@it=\textit\def\PY@tc##1{\textcolor[rgb]{0.24,0.48,0.48}{##1}}}
\@namedef{PY@tok@cs}{\let\PY@it=\textit\def\PY@tc##1{\textcolor[rgb]{0.24,0.48,0.48}{##1}}}

\def\PYZbs{\char`\\}
\def\PYZus{\char`\_}
\def\PYZob{\char`\{}
\def\PYZcb{\char`\}}
\def\PYZca{\char`\^}
\def\PYZam{\char`\&}
\def\PYZlt{\char`\<}
\def\PYZgt{\char`\>}
\def\PYZsh{\char`\#}
\def\PYZpc{\char`\%}
\def\PYZdl{\char`\$}
\def\PYZhy{\char`\-}
\def\PYZsq{\char`\'}
\def\PYZdq{\char`\"}
\def\PYZti{\char`\~}
% for compatibility with earlier versions
\def\PYZat{@}
\def\PYZlb{[}
\def\PYZrb{]}
\makeatother


    % For linebreaks inside Verbatim environment from package fancyvrb.
    \makeatletter
        \newbox\Wrappedcontinuationbox
        \newbox\Wrappedvisiblespacebox
        \newcommand*\Wrappedvisiblespace {\textcolor{red}{\textvisiblespace}}
        \newcommand*\Wrappedcontinuationsymbol {\textcolor{red}{\llap{\tiny$\m@th\hookrightarrow$}}}
        \newcommand*\Wrappedcontinuationindent {3ex }
        \newcommand*\Wrappedafterbreak {\kern\Wrappedcontinuationindent\copy\Wrappedcontinuationbox}
        % Take advantage of the already applied Pygments mark-up to insert
        % potential linebreaks for TeX processing.
        %        {, <, #, %, $, ' and ": go to next line.
        %        _, }, ^, &, >, - and ~: stay at end of broken line.
        % Use of \textquotesingle for straight quote.
        \newcommand*\Wrappedbreaksatspecials {%
            \def\PYGZus{\discretionary{\char`\_}{\Wrappedafterbreak}{\char`\_}}%
            \def\PYGZob{\discretionary{}{\Wrappedafterbreak\char`\{}{\char`\{}}%
            \def\PYGZcb{\discretionary{\char`\}}{\Wrappedafterbreak}{\char`\}}}%
            \def\PYGZca{\discretionary{\char`\^}{\Wrappedafterbreak}{\char`\^}}%
            \def\PYGZam{\discretionary{\char`\&}{\Wrappedafterbreak}{\char`\&}}%
            \def\PYGZlt{\discretionary{}{\Wrappedafterbreak\char`\<}{\char`\<}}%
            \def\PYGZgt{\discretionary{\char`\>}{\Wrappedafterbreak}{\char`\>}}%
            \def\PYGZsh{\discretionary{}{\Wrappedafterbreak\char`\#}{\char`\#}}%
            \def\PYGZpc{\discretionary{}{\Wrappedafterbreak\char`\%}{\char`\%}}%
            \def\PYGZdl{\discretionary{}{\Wrappedafterbreak\char`\$}{\char`\$}}%
            \def\PYGZhy{\discretionary{\char`\-}{\Wrappedafterbreak}{\char`\-}}%
            \def\PYGZsq{\discretionary{}{\Wrappedafterbreak\textquotesingle}{\textquotesingle}}%
            \def\PYGZdq{\discretionary{}{\Wrappedafterbreak\char`\"}{\char`\"}}%
            \def\PYGZti{\discretionary{\char`\~}{\Wrappedafterbreak}{\char`\~}}%
        }
        % Some characters . , ; ? ! / are not pygmentized.
        % This macro makes them "active" and they will insert potential linebreaks
        \newcommand*\Wrappedbreaksatpunct {%
            \lccode`\~`\.\lowercase{\def~}{\discretionary{\hbox{\char`\.}}{\Wrappedafterbreak}{\hbox{\char`\.}}}%
            \lccode`\~`\,\lowercase{\def~}{\discretionary{\hbox{\char`\,}}{\Wrappedafterbreak}{\hbox{\char`\,}}}%
            \lccode`\~`\;\lowercase{\def~}{\discretionary{\hbox{\char`\;}}{\Wrappedafterbreak}{\hbox{\char`\;}}}%
            \lccode`\~`\:\lowercase{\def~}{\discretionary{\hbox{\char`\:}}{\Wrappedafterbreak}{\hbox{\char`\:}}}%
            \lccode`\~`\?\lowercase{\def~}{\discretionary{\hbox{\char`\?}}{\Wrappedafterbreak}{\hbox{\char`\?}}}%
            \lccode`\~`\!\lowercase{\def~}{\discretionary{\hbox{\char`\!}}{\Wrappedafterbreak}{\hbox{\char`\!}}}%
            \lccode`\~`\/\lowercase{\def~}{\discretionary{\hbox{\char`\/}}{\Wrappedafterbreak}{\hbox{\char`\/}}}%
            \catcode`\.\active
            \catcode`\,\active
            \catcode`\;\active
            \catcode`\:\active
            \catcode`\?\active
            \catcode`\!\active
            \catcode`\/\active
            \lccode`\~`\~
        }
    \makeatother

    \let\OriginalVerbatim=\Verbatim
    \makeatletter
    \renewcommand{\Verbatim}[1][1]{%
        %\parskip\z@skip
        \sbox\Wrappedcontinuationbox {\Wrappedcontinuationsymbol}%
        \sbox\Wrappedvisiblespacebox {\FV@SetupFont\Wrappedvisiblespace}%
        \def\FancyVerbFormatLine ##1{\hsize\linewidth
            \vtop{\raggedright\hyphenpenalty\z@\exhyphenpenalty\z@
                \doublehyphendemerits\z@\finalhyphendemerits\z@
                \strut ##1\strut}%
        }%
        % If the linebreak is at a space, the latter will be displayed as visible
        % space at end of first line, and a continuation symbol starts next line.
        % Stretch/shrink are however usually zero for typewriter font.
        \def\FV@Space {%
            \nobreak\hskip\z@ plus\fontdimen3\font minus\fontdimen4\font
            \discretionary{\copy\Wrappedvisiblespacebox}{\Wrappedafterbreak}
            {\kern\fontdimen2\font}%
        }%

        % Allow breaks at special characters using \PYG... macros.
        \Wrappedbreaksatspecials
        % Breaks at punctuation characters . , ; ? ! and / need catcode=\active
        \OriginalVerbatim[#1,codes*=\Wrappedbreaksatpunct]%
    }
    \makeatother

    % Exact colors from NB
    \definecolor{incolor}{HTML}{303F9F}
    \definecolor{outcolor}{HTML}{D84315}
    \definecolor{cellborder}{HTML}{CFCFCF}
    \definecolor{cellbackground}{HTML}{F7F7F7}

    % prompt
    \makeatletter
    \newcommand{\boxspacing}{\kern\kvtcb@left@rule\kern\kvtcb@boxsep}
    \makeatother
    \newcommand{\prompt}[4]{
        {\ttfamily\llap{{\color{#2}[#3]:\hspace{3pt}#4}}\vspace{-\baselineskip}}
    }
    

    
    % Prevent overflowing lines due to hard-to-break entities
    \sloppy
    % Setup hyperref package
    \hypersetup{
      breaklinks=true,  % so long urls are correctly broken across lines
      colorlinks=true,
      urlcolor=urlcolor,
      linkcolor=linkcolor,
      citecolor=citecolor,
      }
    % Slightly bigger margins than the latex defaults
    
    \geometry{verbose,tmargin=1in,bmargin=1in,lmargin=1in,rmargin=1in}
    
    

\begin{document}
    
    \maketitle
    
    

    
    \begin{figure}
\centering
\pandocbounded{\includegraphics[keepaspectratio,alt={Logo portada}]{Informe Final_files/image1.jpg}}
\caption{Logo portada}
\end{figure}

    MODELOS Y SISTEMAS

TRABAJO FINAL

Modelo de contaminante en reactor

Autores:

Ballan, NicolePantozzi, JeremíasOjeda, Augusto

Docentes:

Rubio DianaSoria Paula RominaLa Mura Guillermo

Noviembre 2026

    \begin{center}\rule{0.5\linewidth}{0.5pt}\end{center}

{\def\LTcaptype{none} % do not increment counter
\begin{longtable}[]{@{}
  >{\raggedright\arraybackslash}p{(\linewidth - 2\tabcolsep) * \real{0.5000}}
  >{\raggedleft\arraybackslash}p{(\linewidth - 2\tabcolsep) * \real{0.5000}}@{}}
\toprule\noalign{}
\endhead
\bottomrule\noalign{}
\endlastfoot
\pandocbounded{\includegraphics[keepaspectratio,alt={Logo encabezado}]{Informe Final_files/image1.jpg}}\textbf{Modelo
de contaminante en reactor}Ballan Nicole, Pantozzi Jeremías, Ojeda
Augusto & Noviembre, 2026 \\
\end{longtable}
}

\begin{center}\rule{0.5\linewidth}{0.5pt}\end{center}

    Modelo de contaminante en reactor

Ballan Nicole; Pantozzi Jeremías; Ojeda Augusto

Escuela de Ciencia y Tecnología, Universidad Nacional de San Martín, 25
de Mayo y Francia,1650 San Martín, Buenos Aires, Argentina
correo@autor1.com, jmpantozzi@estudiantes.unsam.edu.ar,
agojeda@estudiantes.unsam.edu.ar

    \textbf{Resumen:} En el presente trabajo se estudia un modelo matemático
utilizado para describir la biodegradación de contaminantes en recursos
hídricos, fundamentado en el principio de conservación de la materia e
incorporando una cinética saturable de tipo Michaelis-Menten. Mediante
herramientas numéricas implementadas en Python, se analiza el sistema no
lineal y su linealización alrededor de un punto de equilibrio, con el
objetivo de determinar la evolución temporal de la concentración, la
masa acumulada y la fuente contaminante. Se calculan los puntos de
equilibrio, se evalúa su estabilidad mediante los autovalores de la
matriz jacobiana y se comparan las simulaciones temporales de ambos
modelos. Además, se obtiene la función de transferencia para estudiar la
respuesta de la concentración frente a perturbaciones externas,
incluyendo las respuestas al impulso y al escalón. Los resultados
indican estabilidad matemática local del modelo, aunque ambos
equilibrios presentan valores negativos en alguna variable, lo que
limita su interpretación física.

Palabras claves: \textbf{\emph{Contaminantes ambientales, Recursos
hídricos, Biodegradación, Python, Linealización, Función de
transferencia.}}

    \subsection{1. Introducción}\label{introducciuxf3n}

La incorporación de contaminantes a los sistemas acuáticos, como
consecuencia de actividades humanas o eventos accidentales, constituye
una problemática ambiental que puede afectar significativamente la
calidad del agua. El estudio de los mecanismos de transporte,
transformación y eliminación de estas sustancias resulta fundamental
para comprender su comportamiento y predecir su evolución temporal. En
este contexto, los modelos matemáticos permiten representar de manera
simplificada los procesos involucrados y analizar la respuesta del
sistema ante diferentes condiciones.

Para abordar esta problemática, se considera un reactor perfectamente
mezclado de volumen constante, en el cual un contaminante es
transportado mediante flujos de entrada y salida, eliminado a través de
procesos de degradación biológica y aportado por una fuente variable en
el tiempo. La formulación se fundamenta en el principio de conservación
de la materia, aplicado al transporte de contaminantes en medios
acuáticos, tomando como referencia los desarrollos de Fischer et
al.~(1979). Este principio se expresa mediante el siguiente balance
general:

\[
\boxed{\text{Acumulación}=\text{Entrada}-\text{Salida}+\text{Generación}-\text{Eliminación}}
\tag{1}
\]

Considerando una concentración uniforme en el reactor, la masa disuelta
del contaminante puede expresarse como \(m_C=VC\), donde \(V\)
representa el volumen del reactor y \(C\) la concentración. Si
inicialmente se supone una eliminación de primer orden, caracterizada
por una constante \(k\), el balance de materia adopta la forma:

\[
\frac{d(VC)}{dt}=QC_{\mathrm{in}}-QC+VS-kVC
\tag{2}
\]

donde \(Q\) representa el caudal volumétrico, \(C_{\mathrm{in}}\) la
concentración del flujo de entrada y \(S\) una fuente interna de
contaminante por unidad de volumen.

Sin embargo, la hipótesis de eliminación de primer orden supone que la
velocidad de degradación aumenta proporcionalmente con la concentración,
comportamiento que puede resultar insuficiente cuando la eliminación
depende de microorganismos cuya capacidad de degradación es limitada.
Por este motivo, se adopta una cinética saturable de tipo
Michaelis-Menten, utilizada para describir procesos de biodegradación
{[}1, 2{]}:

\[
r_d(C)=\frac{K_{\max}C}{K_m+C}
\tag{3}
\]

En esta expresión, \(K_{\max}\) representa la velocidad máxima de
degradación por unidad de volumen y \(K_m\) la constante de
semisaturación. Para concentraciones considerablemente menores que
\(K_m\), la expresión puede aproximarse mediante una cinética de primer
orden, estableciendo la relación \(K_{\max}=kK_m\). De esta manera, se
reemplaza el término de eliminación lineal por una formulación no lineal
que contempla la saturación del proceso biológico.

A partir de estas consideraciones, el comportamiento dinámico del
reactor se describe mediante tres variables de estado: la concentración
del contaminante \(C(t)\), la masa acumulada \(M(t)\) y la intensidad de
la fuente contaminante \(S(t)\). El sistema de ecuaciones diferenciales
propuesto es:

\[
\begin{aligned}
\dot C &= -\frac{Q}{V}C-\frac{K_{\max}C}{K_m+C}+S\\[6pt]
\dot M &= Q(C_{\mathrm{in}}-C)-k_rM\\[6pt]
\dot S &= -\lambda S+u(t)
\end{aligned}
\tag{4}
\]

\begin{itemize}
\item
  La primera ecuación describe la evolución de la concentración
  considerando la salida del contaminante, su degradación biológica y el
  aporte de la fuente interna.
\item
  La segunda representa la evolución de la variable de masa acumulada,
  incluyendo los intercambios asociados al caudal y un término de
  eliminación de primer orden.
\item
  La tercera ecuación caracteriza el decaimiento de la fuente
  contaminante mediante la constante \(\lambda\), incorporando una
  entrada externa \(u(t)\) que permite representar aportes adicionales o
  perturbaciones sobre el sistema.
\end{itemize}

A partir del modelo propuesto, se busca caracterizar su comportamiento
dinámico mediante la determinación de los puntos de equilibrio y el
análisis de su estabilidad. Para ello, se realizarán simulaciones
numéricas bajo diferentes condiciones iniciales y se obtendrá un modelo
linealizado alrededor de un punto de equilibrio. Finalmente, se
compararán las respuestas temporales y los diagramas de fase de ambos
modelos, y se estudiará su respuesta ante perturbaciones externas
mediante el análisis de la función de transferencia.

    \subsection{2. Materiales y Métodos}\label{materiales-y-muxe9todos}

Para el análisis del comportamiento dinámico del sistema se utilizó el
lenguaje de programación \textbf{Python}, trabajando en el entorno
\textbf{Jupyter Notebook}, que permite integrar desarrollos matemáticos,
cálculos numéricos y representaciones gráficas. La implementación
computacional se realizó mediante las siguientes bibliotecas:

\begin{itemize}
\item
  \textbf{NumPy:} se utilizó para el manejo de vectores y matrices, el
  cálculo de autovalores, la manipulación de polinomios y la generación
  de arreglos numéricos necesarios para las simulaciones.
\item
  \textbf{SciPy:} permitió implementar diferentes herramientas de
  análisis numérico. Se empleó \texttt{scipy.optimize.root} para
  determinar los puntos de equilibrio del sistema,
  \texttt{scipy.integrate.solve\_ivp} para resolver numéricamente las
  ecuaciones diferenciales y \texttt{scipy.signal} para obtener la
  función de transferencia y calcular las respuestas al impulso y al
  escalón.
\item
  \textbf{Matplotlib:} se utilizó para representar gráficamente las
  soluciones temporales, los diagramas de fase bidimensionales y
  tridimensionales, los polos y ceros de la función de transferencia y
  las respuestas dinámicas del sistema. Adicionalmente, se empleó
  \texttt{inset\_axes} para ampliar regiones de interés en las
  representaciones gráficas.
\end{itemize}

A partir de estas herramientas se desarrolló el análisis del modelo no
lineal y su posterior linealización, estudiando los puntos de
equilibrio, la estabilidad y la evolución temporal del sistema frente a
distintas condiciones iniciales. Los procedimientos empleados se
describen a continuación.

    \subsubsection{2.1. Parámetros y condiciones de
simulación}\label{paruxe1metros-y-condiciones-de-simulaciuxf3n}

Para el estudio del modelo no lineal presentado en la ecuación (4), se
adoptaron los parámetros y condiciones iniciales que se detallan a
continuación. La formulación se fundamenta en los principios de
transporte y conservación de materia desarrollados por Fischer et
al.~(1979), mientras que los trabajos de Shetty y Shetty (2016) {[}1{]}
y Lv et al.~(2022) {[}2{]} se tomaron como antecedentes para la
descripción de los procesos de biodegradación.

Los valores utilizados permiten caracterizar el comportamiento dinámico
del reactor y realizar las simulaciones numéricas correspondientes.

\textbf{Condiciones iniciales:}

\[
C(0)=0.5\ \mathrm{mg\,L^{-1}}, \qquad
M(0)=1\ \mathrm{mg}, \qquad
S(0)=0.5\ \mathrm{mg\,L^{-1}\,min^{-1}}
\]

\textbf{Tabla 1. Parámetros utilizados en el modelo.}

{\def\LTcaptype{none} % do not increment counter
\begin{longtable}[]{@{}
  >{\centering\arraybackslash}p{(\linewidth - 6\tabcolsep) * \real{0.2632}}
  >{\centering\arraybackslash}p{(\linewidth - 6\tabcolsep) * \real{0.2632}}
  >{\centering\arraybackslash}p{(\linewidth - 6\tabcolsep) * \real{0.2632}}
  >{\raggedright\arraybackslash}p{(\linewidth - 6\tabcolsep) * \real{0.2105}}@{}}
\toprule\noalign{}
\begin{minipage}[b]{\linewidth}\centering
Parámetro
\end{minipage} & \begin{minipage}[b]{\linewidth}\centering
Valor
\end{minipage} & \begin{minipage}[b]{\linewidth}\centering
Unidad
\end{minipage} & \begin{minipage}[b]{\linewidth}\raggedright
Descripción
\end{minipage} \\
\midrule\noalign{}
\endhead
\bottomrule\noalign{}
\endlastfoot
\(Q\) & 1 & \(\mathrm{L\,min^{-1}}\) & Caudal de entrada y salida \\
\(V\) & 10 & \(\mathrm{L}\) & Volumen del reactor \\
\(C_{\mathrm{in}}\) & 1 & \(\mathrm{mg\,L^{-1}}\) & Concentración de
entrada \\
\(K_{\max}\) & 0.644 & \(\mathrm{mg\,L^{-1}\,min^{-1}}\) & Velocidad
máxima de degradación biológica \\
\(K_m\) & 12.88 & \(\mathrm{mg\,L^{-1}}\) & Constante de
semisaturación \\
\(k_r\) & 0.02 & \(\mathrm{min^{-1}}\) & Constante de remoción de
masa \\
\(\lambda\) & 0.1 & \(\mathrm{min^{-1}}\) & Constante de decaimiento de
la fuente \\
\end{longtable}
}

Para la resolución numérica del sistema se consideró un intervalo
temporal de \(0\) a \(100\) minutos, seleccionando \(N=5000\) instantes
de evaluación equiespaciados. Asimismo, se adoptó una entrada externa
constante \(u(t)\), con el propósito de estudiar el comportamiento del
reactor bajo un aporte sostenido a la fuente contaminante.

\textbf{Condiciones de simulación:}

\[
t\in[0,100]\ \mathrm{min}, \qquad N=5000
\]

\[
u(t)=0.2\ \mathrm{mg\,L^{-1}\,min^{-2}}
\]

    \subsubsection{2.2. Puntos de equilibrio y sistema
linealizado}\label{puntos-de-equilibrio-y-sistema-linealizado}

\paragraph{Determinación de los puntos de
equilibrio}\label{determinaciuxf3n-de-los-puntos-de-equilibrio}

Para determinar los puntos de equilibrio del modelo no lineal presentado
en la ecuación (4), se igualaron a cero las tres ecuaciones
diferenciales del sistema, considerando una entrada externa constante
\(u(t)=0.2\).

\[
\begin{aligned}
0 &= -\frac{Q}{V}C^*-\frac{K_{\max}C^*}{K_m+C^*}+S^*\\[4pt]
0 &= Q(C_{\mathrm{in}}-C^*)-k_rM^*\\[4pt]
0 &= -\lambda S^*+u(t)
\end{aligned}
\tag{5}
\]

Donde \(C^*\), \(M^*\) y \(S^*\) representan los valores de las
variables de estado en el equilibrio.

Despejando \(S^*\) de la tercera ecuación y sustituyendo en la primera,
se obtiene una expresión que puede reorganizarse como una ecuación
cuadrática para la concentración de equilibrio:

\[
\frac{Q}{V}(C^*)^2+
\left(\frac{QK_m}{V}+K_{\max}-S^*\right)C^*
-S^*K_m=0
\tag{6}
\]

Por lo tanto, para los parámetros considerados, el sistema presenta dos
puntos de equilibrio matemáticos.

Para determinar numéricamente ambos puntos se utilizó la función
\texttt{root} de la biblioteca \texttt{scipy.optimize}, que permite
encontrar las raíces de un sistema de ecuaciones no lineales. Se
emplearon dos estimaciones iniciales diferentes, con el objetivo de
obtener ambas soluciones.

\paragraph{Linealización del
sistema}\label{linealizaciuxf3n-del-sistema}

Para estudiar el comportamiento del sistema en las proximidades de un
punto de equilibrio, se realizó una linealización mediante una
aproximación de Taylor de primer orden.

Considerando el vector de estado \(\mathbf{x}=[C,M,S]^T\), el sistema
linealizado alrededor de un punto de equilibrio \(\mathbf{x}^*\), para
una entrada externa constante, se expresa como:

\[
\dot{\mathbf{x}}=A(\mathbf{x}-\mathbf{x}^*) + u(t)
\tag{7}
\]

Donde \(A\) es la matriz jacobiana del sistema no lineal evaluada en el
punto de equilibrio seleccionado.

La matriz jacobiana se obtiene calculando las derivadas parciales de
cada ecuación diferencial respecto de las variables de estado:

\[
J(C,M,S)=
\begin{pmatrix}
-\dfrac{Q}{V}-\dfrac{K_{\max}K_m}{(K_m+C)^2} & 0 & 1\\[8pt]
-Q & -k_r & 0\\[4pt]
0 & 0 & -\lambda
\end{pmatrix}
\tag{8}
\]

Para realizar la linealización se seleccionó el primer punto de
equilibrio obtenido numéricamente:

\[
\mathbf{x}^*=
\begin{pmatrix}
16.39352\\
-769.67615\\
2
\end{pmatrix}
\]

Evaluando la matriz jacobiana en este punto, se obtiene:

\[
A=J(\mathbf{x}^*)=
\begin{pmatrix}
-0.10967947 & 0 & 1\\
-1 & -0.02 & 0\\
0 & 0 & -0.1
\end{pmatrix}
\tag{9}
\]

Finalmente, sustituyendo la matriz \(A\) y los valores del equilibrio en
la ecuación (7), se obtiene el sistema linealizado:

\[
\begin{aligned}
\dot C &= -0.10968(C-16.39352)+(S-2)\\[4pt]
\dot M &= -(C-16.39352)-0.02(M+769.67615)\\[4pt]
\dot S &= -0.1(S-2) + u(t)
\end{aligned}
\tag{10}
\]

Este sistema permite aproximar el comportamiento del modelo no lineal en
las proximidades del punto de equilibrio seleccionado. Para analizar su
estabilidad local, se calcularon los autovalores de la matriz jacobiana
\(A\) mediante la función \texttt{numpy.linalg.eigvals}.

    \subsubsection{2.3. Función de
transferencia}\label{funciuxf3n-de-transferencia}

Para analizar la respuesta dinámica del sistema ante perturbaciones
externas, se determinó la función de transferencia del modelo
linealizado. Se consideró como entrada la variación de la fuente externa
\(u(t)\) y como salida la concentración del contaminante \(C(t)\),
debido al interés en estudiar cómo los cambios en el aporte de
contaminantes afectan su concentración dentro del reactor.

Para obtener la función de transferencia se aplica la transformada de
Laplace al sistema linealizado, considerando condiciones iniciales nulas
para las variables de perturbación. La relación entre la entrada y la
salida queda definida por:

\[
H(s)=\frac{\Delta C(s)}{\Delta U(s)}
\tag{12}
\]

A partir de las ecuaciones correspondientes a \(C\) y \(S\), se obtiene:

\[
\begin{aligned}
s\Delta C(s)&=-a\Delta C(s)+\Delta S(s)\\
s\Delta S(s)&=-\lambda\Delta S(s)+\Delta U(s)
\end{aligned}
\]

donde \(a=-A_{11}=0.10967947\) y \(\lambda=-A_{33}=0.1\). Despejando y
combinando ambas expresiones, se obtiene la función de transferencia:

\[
\boxed{
H(s)=\frac{1}{(s+a)(s+\lambda)}
=\frac{1}{s^2+0.20967947s+0.01096795}
}
\tag{13}
\]

Para su implementación numérica en Python, se utilizaron los
coeficientes de la matriz jacobiana evaluada. El numerador se definió
mediante el elemento \(A_{13}\), mientras que el denominador se
construyó multiplicando los polinomios \((s+a)\) y \((s+\lambda)\)
mediante la función \texttt{numpy.polymul}.

Posteriormente, se empleó la función \texttt{numpy.roots} para calcular
las raíces de los polinomios del numerador y denominador,
correspondientes a los ceros y polos de la función de transferencia,
respectivamente.

Finalmente, mediante las funciones \texttt{TransferFunction},
\texttt{impulse} y \texttt{step} de la biblioteca \texttt{scipy.signal},
se construyó la representación del sistema y se calcularon sus
respuestas temporales ante entradas de tipo impulso y escalón, con el
objetivo de caracterizar su comportamiento dinámico frente a
perturbaciones externas.

    \subsection{3. Experimentos numéricos}\label{experimentos-numuxe9ricos}

Con las funciones \texttt{Root} se llego a los puntos de equilibrio, a
partir de los mismos se realizaron las simulaciones temporales de fase y
de respuesta temporal, asi como la linealizacion al rededor del punto
P.E. 1

    \begin{Verbatim}[commandchars=\\\{\}]
=== PUNTOS DE EQUILIBRIO ===
P.E. 1: [  16.39352298 -769.67614889    2.        ]
P.E. 2: [-15.71352298 835.67614889   2.        ]
    \end{Verbatim}

    \begin{Verbatim}[commandchars=\\\{\}]
Jacobiano evaluado en el punto
[[-0.10967947  0.          1.        ]
 [-1.         -0.02        0.        ]
 [ 0.          0.         -0.1       ]]

=== ECUACIONES LINEALIZADAS ===
C\_der = -0.10968*(C - (16.39352)) + (S - (2.00000))
M\_der = -1.00000*(C - (16.39352)) + -0.02000*(M - (-769.67615))
S\_der = -0.10000*(S - (2.00000))

=== Calculo de autovalores ===
λ1 = -0.02, λ2 = -0.10967947376878395, λ3 = -0.1
    \end{Verbatim}

    \begin{center}
    \adjustimage{max size={0.9\linewidth}{0.9\paperheight}}{Informe Final_files/Informe Final_13_0.png}
    \end{center}
    { \hspace*{\fill} \\}
    
    \begin{center}
    \adjustimage{max size={0.9\linewidth}{0.9\paperheight}}{Informe Final_files/Informe Final_14_0.png}
    \end{center}
    { \hspace*{\fill} \\}
    
    \begin{center}
    \adjustimage{max size={0.9\linewidth}{0.9\paperheight}}{Informe Final_files/Informe Final_15_0.png}
    \end{center}
    { \hspace*{\fill} \\}
    
    \begin{Verbatim}[commandchars=\\\{\}]
Ignoring fixed y limits to fulfill fixed data aspect with adjustable data
limits.
C:\textbackslash{}Users\textbackslash{}jerem\textbackslash{}AppData\textbackslash{}Local\textbackslash{}Temp\textbackslash{}ipykernel\_12096\textbackslash{}2671067970.py:81: UserWarning:
This figure includes Axes that are not compatible with tight\_layout, so results
might be incorrect.
  plt.tight\_layout()
Ignoring fixed x limits to fulfill fixed data aspect with adjustable data
limits.
    \end{Verbatim}

    \begin{Verbatim}[commandchars=\\\{\}]

=== POLOS Y CEROS ===
Numerador: [1.]
Denominador: [1.         0.20967947 0.01096795]
Polos: [-0.10967947 -0.1       ]
Ceros: []
    \end{Verbatim}

    \begin{center}
    \adjustimage{max size={0.9\linewidth}{0.9\paperheight}}{Informe Final_files/Informe Final_16_2.png}
    \end{center}
    { \hspace*{\fill} \\}
    
    \begin{center}
    \adjustimage{max size={0.9\linewidth}{0.9\paperheight}}{Informe Final_files/Informe Final_16_3.png}
    \end{center}
    { \hspace*{\fill} \\}
    
    \subsection{4. Análisis y Discusión}\label{anuxe1lisis-y-discusiuxf3n}

\paragraph{4.1. Sistema no linealizado}\label{sistema-no-linealizado}

Los dos puntos de equilibrio obtenidos presentan limitaciones desde el
punto de vista físico, dado que uno de ellos presenta una concentración
negativa y el otro una masa acumulada negativa. Estas magnitudes
representan cantidades físicas que deben satisfacer \(C\geq 0\) y
\(M\geq 0\).

Por otra parte, en la Figura 2 se observa que las trayectorias simuladas
a partir de condiciones iniciales físicamente posibles presentan una
tendencia hacia el primer punto de equilibrio,
\([C=16.393;\ M=-769.676;\ S=2]\).

Aunque dicho equilibrio no es completamente admisible físicamente debido
al valor negativo de la masa, fue seleccionado para continuar con el
análisis matemático del sistema, mientras que el segundo punto fue
descartado por presentar una concentración negativa y por la tendencia
del gráfico mencionado.

\paragraph{4.2. Sistema linealizado}\label{sistema-linealizado}

Al analizar los autovalores de la matriz jacobiana evaluada en el punto
de equilibrio seleccionado, se observa que los tres son reales y
negativos. Esto permite clasificar al equilibrio como un nodo
asintóticamente estable.

Este comportamiento se representa en la Figura 1 mediante los retratos
de fase bidimensionales del sistema linealizado. En ellos se observa
cómo las trayectorias convergen hacia el origen de coordenadas, que
corresponde al punto de equilibrio debido al desplazamiento y
normalización de las variables de estado.

\paragraph{4.3. Comparación de las soluciones
temporales}\label{comparaciuxf3n-de-las-soluciones-temporales}

En las Figuras 2 y 3 se presentan las soluciones temporales del modelo
no lineal y del modelo linealizado, respectivamente.

Se observa que ambos modelos presentan un comportamiento dinámico
similar, con trayectorias que tienden hacia el punto de equilibrio
seleccionado. Este resultado es consistente con el análisis de
estabilidad realizado previamente. Asimismo, en las Figuras 4 y 5 se
representa este comportamiento mediante los diagramas de fase
tridimensionales, considerando las variables de estado \(C\), \(M\) y
\(S\).

Sin embargo, las respuestas no son idénticas. Las principales
diferencias se presentan durante el régimen transitorio, especialmente
en los valores máximos alcanzados por la concentración \(C(t)\) y en la
evolución de la masa acumulada \(M(t)\). Para condiciones iniciales
iguales, estas diferencias se deben a que el modelo linealizado
constituye una aproximación de primer orden del sistema original, cuya
precisión es mayor en las proximidades del punto de equilibrio.

Por otra parte, la variable \(S(t)\) presenta un comportamiento
coincidente en ambos modelos cuando se consideran las mismas condiciones
iniciales, debido a que su ecuación diferencial ya es lineal y no se
modifica mediante el proceso de linealización.

Finalmente, se observa que la concentración y la fuente contaminante se
aproximan a sus respectivos valores de equilibrio dentro del intervalo
simulado, mientras que la masa acumulada presenta una evolución más
lenta. Esto es consistente con el autovalor \(\lambda_1=-0.02\),
asociado a una dinámica de mayor escala temporal.

\paragraph{4.4. Analisís de Función
Transferencia}\label{analisuxeds-de-funciuxf3n-transferencia}

En la Figura 6 se presenta el mapa de polos y ceros de la función de
transferencia obtenida para el modelo linealizado. Se observa que ambos
polos se encuentran en el semiplano izquierdo del plano complejo, lo que
indica que la relación entrada-salida analizada es estable. Este
resultado es consistente con el análisis de estabilidad realizado
previamente.

En la Figura 7 se representa la respuesta al impulso, en la cual se
observa un aumento inicial de la concentración, seguido de una
disminución progresiva hasta aproximarse nuevamente a cero en términos
de perturbación respecto del equilibrio.

Por otra parte, en la Figura 8 se presenta la respuesta al escalón, que
muestra un incremento progresivo de la concentración hasta aproximarse a
un nuevo valor estacionario.

    \subsection{5. Conclusiones}\label{conclusiones}

A partir de los análisis realizados, se logró caracterizar el
comportamiento dinámico del modelo matemático propuesto para el estudio
de los mecanismos de transporte, transformación y eliminación de
sustancias contaminantes en recursos hidricos. El modelo, fudamentado en
el principio de conservación de la materia y adoptando una cinética
saturable del tipo Michaelis-Menten, tomando como referencia los
desarrollos de Fischer (1979), Shetty (2016) y Wang-Liu-Chen-Zhang
(2022), permitió estudiar los puntos de equilibrio, la estabilidad y la
respuesta del sistema frente a diferentes condiciones iniciales y
perturbaciones externas.

Como primer resultado determinante se identifico la limitación fisica
del modelo matematico, dado que ambos puntos de equilibrio presentan
valores negativos de concentración o masa acumulada, magnitudes que no
resultan físicamente admisibles. Si bien se comprobó matemáticamente la
estabilidad asintótica del equilibrio seleccionado y la convergencia de
las trayectorias hacia dicho punto, estos resultados no permiten
asegurar un comportamiento físicamente válido del reactor.

Por otra parte, se verificó que el sistema puede linealizarse alrededor
del punto de equilibrio seleccionado. El modelo linealizado permitió
reproducir cualitativamente el comportamiento del sistema no lineal,
especialmente en las proximidades del equilibrio.

Finalmente, el análisis de la función de transferencia permitió
caracterizar la respuesta del sistema linealizado ante perturbaciones
externas: * La relación entrada-salida analizada es estable, dado que
sus polos se encuentran en el semiplano izquierdo del plano complejo. *
El efecto de una perturbación impulsiva es transitorio y desaparece con
el tiempo. * La respuesta al escalón, muestra un incremento progresivo
de la concentración sin oscilaciones y con un valor final finito.

    \subsection{6. Referencias}\label{referencias}

{[}1{]} Shetty, G. R., y Shetty, V. K., \emph{Pathway identification,
enzyme activity and kinetic study for the biodegradation of phenol by
Nocardia hydrocarbonoxydans NCIM 2386}, Desalination and Water
Treatment, 57(19) (2016), pp.~8789--8801. DOI:
10.1080/19443994.2015.1030700.

{[}2{]} Lv, Y., Wang, L., Liu, X., Chen, B., y Zhang, M.,
\emph{Degradation kinetics of aromatic VOCs polluted wastewater by
functional bacteria at laboratory scale}, Scientific Reports, 12 (2022),
artículo 19053, pp.~1--9. DOI: 10.1038/s41598-022-21356-4.

    \begin{Verbatim}[commandchars=\\\{\}]
[NbConvertApp] Converting notebook MyS\_Trabajo\_Final\_Jupyter.ipynb to latex
[NbConvertApp] Support files will be in Informe Final\_files\textbackslash{}
[NbConvertApp] Making directory Informe Final\_files
[NbConvertApp] Writing 55036 bytes to Informe Final.tex
    \end{Verbatim}


    % Add a bibliography block to the postdoc
    
    
    
\end{document}
